\section{Formulation numérique}
\label{sec-formulation}

\subsection{Discrétisation}

En FDEM 3D, le volume est maillé en tétraèdres linéaires et chaque tétraèdre possède ses quatre nœuds en propre : deux tétraèdres voisins ne partagent pas de nœud, ils sont reliés par un joint cohésif triangulaire à six nœuds posé sur leur face commune (\cle{README.md}, l.~690-700). Le tétraèdre est co-rotationnel : la rotation est extraite du gradient de transformation par trois itérations de Higham et la loi de Hooke s'applique à la déformation de Biot. Cette cinématique est exacte en grande rotation et reste limitée aux petites déformations ; une loi néo-hookéenne, celle de la thèse de Guo~\cite{guo2014thesis} et du code Solidity, est disponible en option (\cref{sec-materiau}). En 2D, les éléments sont des triangles à déformation constante et les joints ont quatre nœuds.

Le maillage peut être une grille découpée en six tétraèdres de Kuhn, une tessellation de Voronoï dont chaque grain est découpé en tétraèdres, ou un maillage Gmsh importé (\cle{mesh = file}). Seul le maillage importé, non structuré, sert aux impacts. La matrice de masse est diagonale : chaque nœud de tétraèdre reçoit $\rho V_0/4$ (\cle{src/Fdem3dSolver.cpp}, l.~2723).

Les joints sont insérés de deux façons. En insertion intrinsèque, tous les joints existent dès l'origine avec une raideur de pénalité finie, ce qui ajoute une souplesse à toutes les faces du maillage. En insertion adaptative~\cite{yan2023adaptive}, aucun joint n'existe à l'origine : les copies d'un même sommet sont liées cinématiquement et se déplacent comme un seul nœud, et un joint n'est créé que lorsque la contrainte sur la face atteint le critère de rupture. Ces deux schémas donnent des résultats différents sur l'impact (\cref{sec-insertion}).

\subsection{Schéma en temps}

L'intégration est explicite, par différences centrées sous la forme demi-pas du saute-mouton : $v \leftarrow v + \Delta t\, f/m$ puis $u \leftarrow u + \Delta t\, v$ (\cle{src/Fdem3dSolver.cpp}, l.~7058-7067). Le schéma est conditionnellement stable. Le pas est fixe par défaut. L'option \cle{dtUpdate = inserted}, ajoutée le 3 octobre 2026 d'après Wu et al.~\cite{wu2024dt}, le rend variable en insertion adaptative : les joints encore liés n'entrent pas dans le calcul du pas initial, la raideur d'un joint s'y ajoute à son insertion, et le pas ne fait que décroître. Le schéma devient alors d'ordre un aux instants de changement de pas (\cle{CHANGELOG.md}, l.~69-71).

\subsection{Pas de temps stable}
\label{sec-dt}

Le pas est calculé une fois à l'initialisation comme le minimum de trois bornes, multiplié par un facteur de sécurité \cle{dtFactor} qui vaut $0{,}15$ en FDEM 3D (\cle{computeStableDt}, \cle{src/Fdem3dSolver.cpp}, l.~3971-4220). La première borne vaut $2\sqrt{m_i/K_i}$ au nœud le plus critique, où $K_i$ cumule la raideur des éléments, celle des joints attachés et une provision pour les contacts. La deuxième est la condition de Courant sur le plus petit diamètre inscrit réel du maillage ; elle remplace depuis le 7 août 2026 le pas nominal de la grille, dont l'usage avait causé l'instabilité de $2{,}4$~MJ. La troisième, $\rho h^2/(4\mu)$, n'est active qu'avec la viscosité de volume. Le journal du calcul imprime le nœud critique et la part des éléments, des joints et du contact.

Les ressorts de pénalité des joints commandent le pas. Sur le maillage de l'impact Kuru le plus fin, le pas vaut $1{,}0$~ns quand la condition de Courant des éléments seuls autorise $4{,}7$~ns, et $98~\%$ de la raideur qui fixe le pas vient des joints (\cle{docs/POINT_rockim_2026-09-13.md}, l.~30-32). Aucune mise à l'échelle des masses n'est employée.

Deux postes ne sont pas comptés dans ce calcul, et l'audit du 13 septembre l'a relevé comme bloquant : la pénalité normale du contact par potentiel et la pénalité nœud-face du contact général n'entrent pas dans le budget du pas en 3D. La raideur tangentielle du potentiel n'y entre qu'avec l'option \cle{dtBudgetTangential = on}, désactivée par défaut ; le code avertit quand elle dépasse la raideur provisionnée, ce qui est le cas sur le banc d'impact de référence (\cref{sec-contact-dt}).

\subsection{Amortissements}

Cinq mécanismes d'amortissement existent ; chacun a sa clé et son défaut (\cref{tab-amortissements}). Leur rôle dans les résultats anciens est documenté. Les calibrations historiques compensaient la perte d'énergie d'un contact quasi plastique (\cle{ROADMAP_rockim.md}, l.~231). Sur un cas mal dimensionné, l'amortissement numérique a absorbé $12~\%$ de l'énergie contre $0{,}55~\%$ pour la fissuration (\cle{GUIDE_rockim.md}, l.~181-189). L'audit de la loi adaptative du 11 septembre a compté six mécanismes dissipatifs armés par défaut sans que la loi constitutive visée les demande, dont l'amortisseur de joint et l'amortissement de Cundall à $0{,}05$ (\cle{AUDIT_loi_adaptative_2026-09-11.md}, §3). Les configurations de reproduction d'articles les désactivent explicitement.

\begin{table}[htbp]
\centering
\caption{Mécanismes d'amortissement de la FDEM 3D et leur réglage par défaut.}
\label{tab-amortissements}
\begin{tabularx}{\linewidth}{@{}X l l@{}}
\toprule
mécanisme & clé & défaut \\
\midrule
Cundall local, $f_d = \alpha |f| \operatorname{sgn}(v)$ par composante & \cle{dampingLocal} & $0{,}05$ \\
amortisseur de joint intact, borné par $m/\Delta t$ & \cle{jointXi} & $0{,}05$ \\
viscosité de volume $2\mu D$ (Yan et al.) & \cle{bulkViscosity} & désactivée \\
amortissement visqueux nodal $-\mu v$ & \cle{dampingViscous} & désactivé \\
amortisseur du contact d'outil et du contact par pénalité & \cle{contactXi}, \cle{gcXi} & $0{,}05$ et $0{,}8$ \\
\bottomrule
\end{tabularx}
\end{table}

\subsection{Frontières absorbantes}

Les faces planes du bloc de roche peuvent porter les amortisseurs de Lysmer et Kuhlemeyer~\cite{lysmer1969}, $\rho c_p v_n$ et $\rho c_s v_t$, répartis sur l'aire tributaire et appliqués de façon implicite dans la mise à jour des vitesses, $v \leftarrow (v + \Delta t F/m)/(1 + \Delta t\, c/m)$, ce qui les laisse sans effet sur le pas. Des ressorts de rappel de raideur $G/R$ en normal et $G/(2R)$ en tangentiel les complètent (\cle{src/Fdem3dSolver.cpp}, l.~3884-3893). La clé \cle{absorbing} vaut \cle{none} (défaut), \cle{sides} ou \cle{all} ; avec \cle{all}, la base du bloc n'est plus encastrée. L'absorption n'est qu'approchée en incidence oblique, et une frontière absorbante ne remplace pas la roche absente autour d'une pointe de fissure (\cle{README.md}, l.~108-117).

\subsection{Unités, configuration et sorties}

rockim travaille en unités SI (m, s, Pa, kg), avec le point décimal obligatoire ; en 2D, les forces sont par mètre d'épaisseur. Cette convention diffère de celle des calculs Abaqus de la thèse (mm, t, s, MPa) ; le script \cle{tools/export_abaqus.py} convertit.

Un calcul se décrit dans un fichier texte \cle{.cfg} de lignes \cle{cle = valeur}. Les propriétés d'un corps ou d'une phase se surchargent par des clés préfixées (\cle{phase.<nom>.<prop>}, \cle{groupVel.<corps>}). Depuis le 5 septembre 2026, une clé inconnue arrête le calcul avec un message qui propose le nom le plus proche, et chaque calcul écrit \cle{config_effective.cfg} avec toutes les valeurs effectivement utilisées, défauts compris. Ce refus est la réponse à un mode de panne mesuré : des clés lues sans effet. Le brossage des fragments, par exemple, était inerte dans sept configurations, de sorte que les calculs Kuru, de pulvérisation et St Anne antérieurs tournaient sans lui (\cle{CHANGELOG.md}, l.~1611-1614).

Un calcul écrit trois sorties. Le fichier \cle{history.csv} contient environ 2\,000 lignes par calcul : temps, cinématique des corps suivis, forces de contact par paire de corps et six travaux par famille de forces. Des trames VTU par élément et par joint portent l'endommagement, l'instant et le mode de rupture de chaque joint. Un résumé sur la sortie standard donne le bilan d'énergie, le décompte des ruptures en traction et en cisaillement, les statistiques de contact et le temps de calcul. Le bilan d'énergie est décrit au \cref{sec-B4}.
